\addcontentsline{toc}{subsubsection}{Solution 37. The Crossed Contraction Lock}

\begin{dlsummary}[title={\bfseries Solution 37. The Crossed Contraction Lock},breakable]

\begin{proof}
The only solution is
\[
\boxed{f(x)=x\qquad(x>0).}
\]

Let $c=f(1)$. Substituting $y=1$ in the main inequality gives
\[
f(cx)+f(f(x))\le2x.
\]
Since $f(f(x))\ge x$, we obtain
\[
f(cx)\le x,
\]
or, equivalently,
\[
f(t)\le\frac tc\qquad(t>0). \tag{1}
\]

Applying (1) to $t=f(x)$ yields
\[
x\le f(f(x))\le\frac{f(x)}c,
\]
and therefore
\[
f(x)\ge cx. \tag{2}
\]

Taking $x=1$ in (1), we have $f(c)\le1$. On the other hand,
$f(c)=f(f(1))\ge1$, so $f(c)=1$. Now set $x=c$ in the main
inequality. Using $f(c)=1$ gives
\[
f\bigl(cf(y)\bigr)+f(y)\le2cy.
\]
By (2), $f(cf(y))\ge c^2f(y)$; hence
\[
(1+c^2)f(y)\le2cy,
\]
so
\[
f(y)\le ky,
\qquad k=\frac{2c}{1+c^2}\le1. \tag{3}
\]

Applying (3) twice gives
\[
f(f(x))\le k^2x.
\]
Together with $f(f(x))\ge x$, this forces $k=1$. Equality in
$2c\le1+c^2$ then gives $c=1$. Consequently, (2) and (3) become
\[
x\le f(x)\le x,
\]
so $f(x)=x$ for every $x>0$.

Finally, $f(x)=x$ satisfies the iteration condition and makes the
main inequality an equality.
\end{proof}
\end{dlsummary}

\begin{takeaways}
\begin{itemize}
\item The substitution $y=1$ converts the two-variable inequality into a
linear upper bound involving the constant $c=f(1)$. The iteration
condition turns that upper bound into a lower bound. A second special
substitution, $x=c$, produces the sharp contraction factor
$2c/(1+c^2)$, and iterating the contraction closes the argument.
\end{itemize}
\end{takeaways}